\documentclass[conference]{IEEEtran}
\\usepackage{fontspec}\\usepackage[english]{babel}\\setmainfont{DejaVu Serif}
\usepackage{amsmath,amssymb,algorithm,algpseudocode,listings,hyperref}
\usepackage{enumitem}\hypersetup{hidelinks}\lstset{language=Python,basicstyle=\footnotesize\ttfamily,breaklines=true,frame=single,columns=fullflexible}
\title{Minimax and Alpha--Beta Pruning: Detailed Analysis and Tic-Tac-Toe Application}
\author{\IEEEauthorblockN{Student Author}\IEEEauthorblockA{Department of Artificial Intelligence\\University or Institution}}
\begin{document}\maketitle
\begin{abstract}This paper explains Minimax and Alpha--Beta pruning for two-player zero-sum games. It defines the game state and value function, gives recursive algorithms and correctness intuition, analyzes bounds, complexity, and move ordering, and applies both methods to Tic-Tac-Toe. The paper also provides an executable Python implementation and describes the exact search process, including terminal tests, value backup, and pruning.\end{abstract}
\begin{IEEEkeywords}Minimax, alpha-beta pruning, game tree, Tic-Tac-Toe, adversarial search\end{IEEEkeywords}
\section{Introduction}In an adversarial game, a move is good only after considering the opponent's best reply. Minimax models this reasoning by recursively alternating MAX and MIN choices. Alpha--Beta pruning keeps the same answer while skipping branches that cannot affect the choice. Tic-Tac-Toe is small enough to search to terminal states, so no heuristic or depth limit is needed.
\section{Game Model}A state is $s=(b,p)$, where $b\in\{X,O,\cdot\}^9$ is the board and $p$ is the player to move. $A(s)$ is the set of empty cells. $T(s,a)$ places $p$ in cell $a$ and changes the player. A state is terminal when a player has three aligned marks or the board is full. From the root player's perspective, utility is $+1$ for a win, $0$ for a draw, and $-1$ for a loss. We assume perfect information, alternating turns, no chance events, and zero-sum rewards.
\section{Minimax}Let $V(s)$ be the value under optimal play. For a terminal state, $V(s)=U(s)$. For a MAX state,
\[V(s)=\max_{a\in A(s)} V(T(s,a)),\]
while for a MIN state,
\[V(s)=\min_{a\in A(s)} V(T(s,a)).\]
The root action is $a^*=\arg\max_a V(T(s,a))$.\begin{algorithm}[h]\caption{Recursive Minimax}\begin{algorithmic}[1]\Require state $s$, root player $r$\If{$s$ is terminal}\State \Return $U(s,r)$\EndIf\If{$player(s)=r$}\State \Return $\max_{a\in A(s)}\mathrm{Minimax}(T(s,a),r)\Else\State \Return $\min_{a\in A(s)}\mathrm{Minimax}(T(s,a),r)$\EndIf\end{algorithmic}\end{algorithm}
\subsection{Correctness}At a terminal node the value is defined. If all child values are correct, MAX selects the largest child and MIN selects the smallest child. Induction on the remaining depth therefore proves that the root value equals the result of optimal play. The empty Tic-Tac-Toe board has value $0$: perfect play produces a draw.
\section{Alpha--Beta Pruning}During depth-first search, $\alpha$ is the best value MAX has secured on the current path, while $\beta$ is the best value MIN has secured. At MAX, update $\alpha\leftarrow\max(\alpha,v)$; at MIN, update $\beta\leftarrow\min(\beta,v)$. When $\alpha\ge\beta$, stop exploring that node. An ancestor already has an alternative at least as good, so the remaining children cannot change the minimax result.\begin{algorithm}[h]\caption{Alpha--Beta Search}\begin{algorithmic}[1]\Require $s,\alpha=-\infty,\beta=+\infty$\If{$s$ is terminal}\State \Return $U(s)$\EndIf\If{MAX}\State $v\gets-\infty$\For{$a\in A(s)$}\State $v\gets\max(v,AB(T(s,a),\alpha,\beta))$; $\alpha\gets\max(\alpha,v)$\If{$\alpha\ge\beta$}\State \textbf{break}\EndIf\EndFor\Else\State $v\gets+\infty$\For{$a\in A(s)$}\State $v\gets\min(v,AB(T(s,a),\alpha,\beta))$; $\beta\gets\min(\beta,v)$\If{$\alpha\ge\beta$}\State \textbf{break}\EndIf\EndFor\EndIf\State \Return $v$\end{algorithmic}\end{algorithm}
\section{Complexity and Ordering}With branching factor $b$ and depth $d$, plain Minimax visits $O(b^d)$ nodes and uses $O(d)$ DFS stack space. Alpha--Beta has the same worst case, but good ordering can reduce the effective work to $O(b^{d/2})$. A practical ordering for Tic-Tac-Toe is center, corners, then edges. The supplied code uses index order to make the trace easy to reproduce; the result remains optimal.
\section{The Search on Tic-Tac-Toe}Each recursive call performs six steps: (1) test whether someone won or the board is full; (2) enumerate empty cells; (3) place the current mark; (4) recurse with the other player; (5) back up the maximum or minimum child value; and (6) update $\alpha$ or $\beta$ and possibly cut. Starting from the empty board, X tries a cell, O tries each reply, and this continues until a terminal state. If a MIN node has already found a value no greater than MAX's $\alpha$, then $\alpha\ge\beta$ and the rest of that node is pruned. All root moves have value $0$ under perfect play.
\subsection{Auditable Search Trace}The program counts visited `nodes` and `cutoffs`. With `--show-search`, it prints every root candidate and its value. It checks all eight winning lines and searches to terminal states, so the output is an exact solution rather than a depth-limited estimate.\section{Python Implementation}\lstinputlisting{ttt_minimax.py}
\section{Conclusion}Minimax gives an optimal adversarial decision rule. Alpha--Beta returns the same decision while reducing work through bounds and ordering. Larger games require depth limits, evaluation functions, transposition tables, and stronger move ordering.\begin{thebibliography}{9}\bibitem{russell}S. Russell and P. Norvig, \emph{Artificial Intelligence: A Modern Approach}, 4th ed., Pearson, 2021.\bibitem{knuth}D. E. Knuth and R. W. Moore, ``An analysis of alpha-beta pruning,'' \emph{Artificial Intelligence}, vol. 6, pp. 293--326, 1975.\bibitem{shannon}C. Shannon, ``Programming a computer for playing chess,'' \emph{Philosophical Magazine}, vol. 41, pp. 256--275, 1950.\end{thebibliography}\end{document}
